\chapter{ENGINEERING CONSIDERATIONS, CHALLENGES
AND REMEDIES}\label{anotherchapter}
\vspace{-3em}

This chapter discusses the ethical, and environmental considerations and social constraints taken into consideration throughout this research work. The engineering activities involved in this thesis have also been elaborated. 

\section{Design Constraints}

\subsection{Ethical Considerations}

...



 \subsection{Sustainability in Environmental and Societal Context}

.....

% \textbf{The proposed system has the potential to foster both environmental and societal sustainability. If successful, it can promote the creation and preservation of green spaces, leading to benefits such as improved air quality and ecosystem health. Furthermore, by facilitating sustainable urban planning, the research aims to enhance societal well-being. Access to green spaces can positively impact physical and mental health, community cohesion, and overall quality of life for residents. However, responsible implementation is crucial to ensure equitable distribution of benefits and prevent unintended social consequences. Ultimately, this research envisions a future where environmental and societal sustainability are interconnected, creating mutual benefits for communities and the planet.}

\section{Complex Engineering Problem}
 \subsection{Complex Problem Solving}
 Complex Engineering Problems addressed by this thesis have been justified in Table \ref{table_OBE1}. %Table \ref{table_OBE1}.
\begin{table}[!h]
    \centering
    \caption{Complex Engineering Problem (CEP)}
    \begin{tabular}{p{.27\linewidth} p{.67\linewidth}}
        \hline
        Name of Complex \newline Engineering Attribute & Justification \\ \hline
                  \textbf{P1 \newline   Depth of knowledge required}   &  For this thesis, in-depth knowledge of Artificial Intelligence, Deep Learning, Machine Learning and Digital Image Processing were required. \\ [0.75em]
                  \textbf{P3 \newline   Depth of analysis \newline required } & Several segmentation models and models suitable for time-series data were explored to propose a suitable model for both green space detection and prediction.  \\ [0.75em]
                  \textbf{P4 \newline  Familiarity of \newline issues } & It was required to gather knowledge of different techniques and measures on how to utilise satellite imagery to measure green space, such as remote sensing and GVI calculation.  \\ [0.75em]
\textbf{P7 \newline   Interdependence} & This thesis work was done in two phases - green space detection and prediction. The prediction phase required segmented images that were achieved from the first phase. Additionally, the first and foremost task of this thesis was data acquisition, without which it is not possible to conduct the two phases.\\
        \hline
    \end{tabular}
    \label{table_OBE1}
\end{table}

\subsection{Engineering Activities}
 Complex Engineering Activities addressed by this thesis have been justified in Table \ref{table_OBE2}.                            
 \begin{table}[ht]
    \centering
    \caption{Complex Engineering Activities (CEA)}
    \begin{tabular}{p{.27\linewidth} p{.67\linewidth}}
        \hline
        Name of Complex \newline Engineering Activity & Justification \\ \hline
\textbf{A1 \newline Range of Resources}   &  Google Earth Pro and Kaggle were the resources for the datasets, Google Drive was utilized to store all data and necessary codes and Google Collaboartary environment was used to facilitate machine learning models. \\ [0.75em]
\textbf{A3 \newline  Innovation } & After an extensive literature review, it was found that no work has been done that uses ensemble learning techniques to segment green space from images. Moreover, no work was found that used a time-series data analysis to predict change in green space. This thesis work is the first research that segments green spaces from images using an ensemble learning approach and then predicts green space. \\ [0.75em]
\textbf{A4 \newline  Consequences for \newline Society and the \newline Environment} &  Environment specialists, policymakers can use this proposed system to take necessary actions to restore the greenery of concerned areas.
Urban planners can also make plans based on this information, taking into account areas that are most at risk of deforestation.
\\
        \hline
    \end{tabular}
    \label{table_OBE2}
\end{table}